% Mapping section: reuse Eloundou-only material from slides/mapping/.
% Core frames only in the main flow; the worked examples, the task-level
% scoring tables, the many-to-one example and the O*NET intro live in the
% Backup section (90_backup.tex).

% Brief intro: why a mapping is needed
\begin{frame}{Mapping from US to Norwegian occupation codes}
\begin{itemize}\setlength{\itemsep}{8pt}
  \item Eloundou et al.'s exposure measure is built on US occupation codes (SOC).
  \item Norwegian register data use STYRK-08.
  \item STYRK-08 is based on ISCO-08; overlapping 4-digit unit groups are code-compatible after a small set of documented Norwegian-title checks.
  \item Bridge: SOC $\to$ ISCO-08 via US Bureau of Labor Statistics (BLS), then match overlapping STYRK-08 codes with manual SOC-source corrections for 2267 and 2269.
\end{itemize}
\end{frame}

% Eloundou concept, chain, pharmacists worked example
% Eloundou et al. (2024): task-level LLM exposure -- core frames
% (split out of eloundou.tex so the presentation deck can input only these)

\begin{frame}{What Eloundou et al.\ measure}
\begin{itemize}\setlength{\itemsep}{6pt}
  \item Each O*NET task rated by humans and GPT-4 on three categories:
  \begin{itemize}
    \item[$E_0$] no exposure
    \item[$E_1$] LLM alone reduces task time by $\geq 50\%$
    \item[$E_2$] LLM with complementary software reduces task time by $\geq 50\%$
  \end{itemize}
  \item Task-level score:\quad
        $\beta_{\text{task}} = 0$ if $E_0$,\;\; $1$ if $E_1$,\;\; $0.5$ if $E_2$
  \item Occupation score: $\beta_{\text{occ}}$ = mean of $\beta_{\text{task}}$ with O*NET Core tasks weighted 2, Supplemental tasks 1
  \item We use the published occupation-level file (\texttt{occ\_level.csv}); it differs slightly from unweighted task means where an occupation has Supplemental tasks
\end{itemize}
\end{frame}

% ---------------------------------------------------------------
\begin{frame}{The code chain}
\centering
\begin{tikzpicture}[node distance=1.0cm and 1.1cm]
  \node[codebox=cTASK] (task) {\cbox{cTASK}{O*NET task}{8823, \dots}{rated $E_0/E_1/E_2$}};
  \node[codebox=cSOC18, right=of task] (soc18) {\cbox{cSOC18}{SOC 2018}{29-1051}{Pharmacists}};
  \node[codebox=cSOC10, right=of soc18] (soc10) {\cbox{cSOC10}{SOC 2010}{29-1051}{Pharmacists}};
  \node[codebox=cSTYRK, right=of soc10] (isco) {\cbox{cSTYRK}{ISCO-08}{2262}{Pharmacists}};

  \draw[chain] (task) -- node[arrowlabel]{average} (soc18);
  \draw[chain] (soc18) -- node[arrowlabel]{BLS 2017} (soc10);
  \draw[chain] (soc10) -- node[arrowlabel]{BLS 2012} (isco);
\end{tikzpicture}

\vspace{1em}
Overlapping 4-digit unit groups in STYRK-08 are code-compatible with ISCO-08 (SSB Notater 17/2011). STYRK 2262 is Farmasøyter.
\end{frame}

% ---------------------------------------------------------------
\begin{frame}{Worked example: Pharmacists (one-to-one mapping)}
\centering
\begin{tikzpicture}[node distance=1.0cm and 1.1cm]
  \node[codebox=cTASK] (task) {\cbox{cTASK}{O*NET task}{21 tasks}{$E_0/E_1/E_2$ ratings}};
  \node[codebox=cSOC18, right=of task] (soc18) {\cbox{cSOC18}{SOC 2018}{29-1051}{Pharmacists}};
  \node[codebox=cSOC10, right=of soc18] (soc10) {\cbox{cSOC10}{SOC 2010}{29-1051}{Pharmacists}};
  \node[codebox=cSTYRK, right=of soc10] (isco) {\cbox{cSTYRK}{ISCO-08}{2262}{Pharmacists}};

  \draw[chain] (task) -- node[arrowlabel]{average} (soc18);
  \draw[chain] (soc18) -- node[arrowlabel]{1$\to$1} (soc10);
  \draw[chain] (soc10) -- node[arrowlabel]{1$\to$1} (isco);
\end{tikzpicture}

\vspace{1em}
\begin{itemize}\setlength{\itemsep}{4pt}
  \item Pharmacists carry the same 6-digit SOC in both vintages: 29-1051
  \item BLS crosswalk sends 29-1051 to a single ISCO-08 code (2262)
  \item Norwegian STYRK 2262 is Farmasøyter, no Norway-specific re-coding
\end{itemize}
\end{frame}

% ---------------------------------------------------------------
\begin{frame}{Task-level scoring for Pharmacists (SOC 29-1051, all 21 tasks)}
\centering\scriptsize
\setlength{\tabcolsep}{4pt}
\renewcommand{\arraystretch}{0.95}
\begin{tabular}{p{5.5cm}cc @{\hspace{0.7cm}} p{5.5cm}cc}
\toprule
O*NET task & $E$ & $\beta$ & O*NET task & $E$ & $\beta$ \\
\midrule
\rO{Review prescriptions for accuracy}{$E_2$}{0.5}             & \rO{Refer patients to other professionals}{$E_2$}{0.5} \\
Assess identity, strength, purity of meds     & $E_0$ & 0.0 & \rO{Work in hospitals/clinics or for HMOs}{$E_2$}{0.5} \\
\rO{Provide info on drug interactions}{$E_2$}{0.5}             & \rR{Update or troubleshoot pharmacy databases}{$E_1$}{1.0} \\
\rO{Analyze prescribing trends}{$E_2$}{0.5}                    & \rO{Manage pharmacy operations, hire, supervise}{$E_2$}{0.5} \\
\rO{Maintain pharmacy files and patient profiles}{$E_2$}{0.5}  & Prepare sterile solutions                           & $E_0$ & 0.0 \\
\rO{Collaborate with other health professionals}{$E_2$}{0.5}   & Offer health promotion / prevention activities      & $E_0$ & 0.0 \\
Plan procedures for mixing and labeling       & $E_0$ & 0.0 & Assay radiopharmaceuticals                          & $E_0$ & 0.0 \\
\rO{Order and purchase pharmaceutical supplies}{$E_2$}{0.5}    & \rR{Publish educational information}{$E_1$}{1.0} \\
Compound and dispense medications             & $E_0$ & 0.0 & & & \\
\rO{Contact insurance companies on billing}{$E_2$}{0.5}        & \multicolumn{3}{l}{\textbf{Counts:} 6 $E_0$ \quad 3 $E_1$ \quad 12 $E_2$} \\
\rO{Advise customers on medication brands}{$E_2$}{0.5}         & & & \\
\rR{Teach pharmacy students serving as interns}{$E_1$}{1.0}    & \multicolumn{3}{l}{\textbf{Occupation $\beta$ (mean of 21):} \;\; \textbf{0.4286 (Q4)}} \\
\rO{Provide services for diabetes, asthma, etc.}{$E_2$}{0.5}   & & & \\
\bottomrule
\end{tabular}
\end{frame}


% Distribution across Norwegian occupations (employment-weighted)
% Eloundou β distribution across STYRK occupations (unweighted + employment-weighted)

\begin{frame}{Distribution of Eloundou et al.\ $\beta$ across Norwegian occupations}
\centering
\includegraphics[width=\linewidth,height=5cm,keepaspectratio]{eloundou_distribution.pdf}

\vspace{0.3em}
\footnotesize
Left: one observation per STYRK-08 code (397 codes; mean 0.34). Right: the same distribution weighted by December 2025 employment (3.18M workers; mean 0.33); the largest single occupations are labelled. Large occupations sit across the whole exposure range, from health care workers (0.11) to policy administration professionals (0.57).
\end{frame}


% ---------------------------------------------------------------
\begin{frame}{Five largest occupations per exposure quintile (Apr 2026)}
\hypertarget{after-mapping}{}
\begin{columns}[T,onlytextwidth]
\begin{column}{0.50\textwidth}
\scriptsize
\setlength{\tabcolsep}{3pt}
\renewcommand{\arraystretch}{0.92}
\begin{tabular}{@{}p{4.6cm} r r@{}}
\toprule
Occupation & Score & Emp. \\
\midrule
\multicolumn{3}{@{}l}{\textit{Q1 (least exposed)}} \\
Cleaners (offices, hotels, \ldots) & 0.03 & 40k \\
Building electricians & 0.13 & 26k \\
Food machine operators & 0.09 & 23k \\
Motor vehicle mechanics & 0.09 & 17k \\
Plumbers and pipe fitters & 0.06 & 17k \\
\addlinespace
\multicolumn{3}{@{}l}{\textit{Q2}} \\
Carpenters and joiners & 0.14 & 35k \\
Child care workers & 0.25 & 32k \\
Waiters & 0.16 & 20k \\
Machinery mechanics and repairers & 0.13 & 16k \\
Cooks & 0.19 & 16k \\
\addlinespace
\multicolumn{3}{@{}l}{\textit{Q3}} \\
Shop sales assistants & 0.34 & 115k \\
Stock clerks & 0.33 & 29k \\
Heavy truck and lorry drivers & 0.31 & 21k \\
Science technicians (n.e.c.) & 0.36 & 20k \\
Car, taxi and van drivers & 0.28 & 14k \\
\bottomrule
\end{tabular}
\end{column}
\begin{column}{0.50\textwidth}
\scriptsize
\setlength{\tabcolsep}{3pt}
\renewcommand{\arraystretch}{0.92}
\begin{tabular}{@{}p{4.6cm} r r@{}}
\toprule
Occupation & Score & Emp. \\
\midrule
\multicolumn{3}{@{}l}{\textit{Q4}} \\
Retail/wholesale trade managers & 0.48 & 30k \\
Managing directors / CEOs & 0.49 & 28k \\
Civil engineers & 0.48 & 19k \\
Sales and marketing managers & 0.48 & 13k \\
Construction managers & 0.44 & 12k \\
\addlinespace
\multicolumn{3}{@{}l}{\textit{Q5 (most exposed)}} \\
Commercial sales representatives & 0.57 & 37k \\
General office clerks & 0.60 & 33k \\
Systems analysts & 0.63 & 22k \\
Software developers (n.e.c.) & 0.77 & 18k \\
Accounting associate profs.\ & 0.80 & 15k \\
\bottomrule
\end{tabular}

\vspace{0.6em}
\footnotesize Employment in thousands. Private sector, ages 21--60, ranked within quintile by April 2026 headcount.

\vspace{0.3em}
\hfill\hyperlink{worked-examples}{\beamerbutton{More worked examples}}
\end{column}
\end{columns}
\end{frame}
